\chapter{Global sign entry for the adjacent contact family}
\label{component:pe31b9ef5ec43}
\noindent\textbf{Component }\path{pe31b9ef5ec43}\par
{\small\noindent Source: \path{CORE/prior/inherited/C1_GLOBAL_SIGN_ENTRY.md}\par}
{\small\noindent SHA256: \path{e31b9ef5ec430cd15faef01eb6b7a33a604ee73f342def5fc245a7aa4a020ccc}\par}
{\small\noindent Inventory use: Retained exact cofactor, selected minor, Cramer identity, seed and C1 sign proof; analytical foundational calls replaced by SE1-SE10\par}
\medskip
\hypertarget{global-sign-entry-for-the-adjacent-contact-family}{%
\section{Global sign entry for the adjacent contact family}\label{pe31b9ef5ec43:global-sign-entry-for-the-adjacent-contact-family}}

\hypertarget{statement-and-exact-scope}{%
\subsection{Statement and exact scope}\label{pe31b9ef5ec43:statement-and-exact-scope}}

Use the canonical full parent in coefficient order
\(c=(d,\lambda,e_1,e_2,e_3,e_4,e_5,\alpha,\beta)\).
On the adjacent face \(\lambda=e_1=0\), assume d\textgreater0 and initially that the other six event
coefficients are strictly positive. Normalize \(e_3=2\). At a physical
noncommuting zero take the angle gauges \(u,v>0\) and normalize the last
kernel coordinate to one:
\[
\psi=(A,u(\eta-1),uY,1)^T.
\]
Then
\[
\boxed{A>0,\qquad 0<Y<A,\qquad 0<\eta<1.}
\tag{1}
\]
These are actual ordered-field inequalities, with no residue or rate
assumptions. The assertion holds over every real closed extension of the
reals. It applies to every such boundary contact, without an extremality
assumption. Section 6 extends the result to the exact three-consecutive-omission
boundary. This is a sign theorem, not an SOS certificate for the full family.

The proof uses the established full-parent determinant, strict curvature,
and contact sufficiency theorems. Its only new continuity argument is
spelled out below; connectedness is used on a specified positive parameter
orthant, not assumed for an unspecified contact locus.

\hypertarget{every-positive-event-tuple-has-one-analytic-maximal-signed-contact}{%
\subsection{1. Every positive event tuple has one analytic maximal-signed contact}\label{pe31b9ef5ec43:every-positive-event-tuple-has-one-analytic-maximal-signed-contact}}

Put \(w=(e_2,e_4,e_5,\alpha,\beta)\in\mathbb R_{>0}^5\), retaining
\(e_3=2\), and define
\[
d_*(w)=\max\{d\ge0:(d,0,0,e_2,2,e_4,e_5,\alpha,\beta)\in C\}.
\tag{2}
\]
This maximum exists, is strictly positive, and is finite. Closedness and
convexity give a closed interval of admissible nonnegative d.~At d=0 all
outer determinant coefficients are nonnegative, and the central coefficient
is strictly positive: its nonnegative expansion contains
\(16e_2e_3e_4e_5\). For sufficiently small d\textgreater0 the central coefficient
remains positive and all outer coefficients remain nonnegative, proving
physicality by full-parent inertia. Testing a fixed strictly CHSH-violating
realization bounds d above.

For these six positive events the determinant support is a fixed
two-dimensional support with zero in its interior. The logarithmic Laurent
determinant therefore is strictly convex and coercive. At d=d\_* its minimum
must be zero: a strictly positive minimum is stable under small increases
of d, by the uniform supported-perturbation estimate in the routing proof.
Thus the minimum supplies unique squared contact angles \(U=u^2,V=v^2>0\).
The parent has three positive eigenvalues, so its kernel is simple. The
established packet argument shows that its last coordinate is nonzero.
Consequently the chosen positive angles and normalized kernel are unique.

At this contact the core expectation is strictly negative. Indeed,
\[
0=d_*\,\psi^TH\psi+\sum_i w_i\,\psi^TE_i\psi.
\]
If all retained expectations vanished, all retained event matrices would
annihilate the state and the kernel equation would give H psi=0. The core
is invertible at a finite noncommuting pair, a contradiction. At least one
expectation is positive, and d\_*\textgreater0, giving \(\psi^TH\psi<0\).

Let z consist of the three affine state coordinates and two positive
angles, and let f\(d,w,z\) be the full polynomial parent slack. Its complete
five-variable contact Hessian is positive definite. The equations
\[
\nabla_z f=0,\qquad f=0
\]
have Jacobian with respect to \(z,d\)
\[
\begin{pmatrix} \nabla_z^2 f&\partial_d\nabla_z f\\
0&\partial_d f\end{pmatrix},
\]
which is invertible because the upper block is positive definite and
\(\partial_d f=\psi^TH\psi<0\). The Nash implicit function theorem gives
a local analytic branch of d and z. Stationary-contact sufficiency makes
that branch physical; its negative core expectation prevents increasing d
at the fixed test. Thus it equals d\_*\(w\), with the unique chosen contact.
These branches agree on overlaps, giving analytic functions on the entire
connected orthant \(\mathbb R_{>0}^5\).

\hypertarget{a0-from-the-negative-chsh-cone}{%
\subsection{2. A\textgreater0 from the negative CHSH cone}\label{pe31b9ef5ec43:a0-from-the-negative-chsh-cone}}

Write an arbitrary real state as \(A,X,Y_s,Z\). Set
\(h=1+v^2+u^2v^2\). The X,Y\_s block of the raw core H is
\[
2\begin{pmatrix}h&-uv\\-uv&h\end{pmatrix}\succ0.
\]
Its Schur complement on A,Z is
\[
\frac{2(1+u^2)(1+v^2)}{h^2-u^2v^2}
\begin{pmatrix}u^2v^2&-uvh\\-uvh&u^2v^2\end{pmatrix}.
\tag{3}
\]
For AZ\textless=0 its quadratic form is nonnegative. Completing the X,Y\_s square
then gives psi\^{}T H psi\textgreater=0. Every contact above has negative core
expectation, so AZ\textgreater0. With Z=1, this proves A\textgreater0.

Every physical noncommuting contact with this fixed event tuple is on the
maximal graph (2): its negative core expectation means that increasing d
at the fixed test immediately produces a negative slack.

\hypertarget{eta0-from-a-positive-cofactor-identity}{%
\subsection{3. eta\textgreater0 from a positive cofactor identity}\label{pe31b9ef5ec43:eta0-from-a-positive-cofactor-identity}}

Let B=adj\(P\). At contact B=gamma psi psi\^{}T for gamma\textgreater0, and hence
\[
u B_{44}+B_{24}=\gamma u\eta.
\]
Here indices are one-based. Direct full-parent expansion gives
\[
u B_{44}+B_{24}=16uv^2(1+u^2)^2(1+v^2)\,F_\eta,
\tag{4}
\]
where, writing Dv=1+v\^{}2,
\[
\begin{aligned}
F_\eta={}&(d+2e_5)\big[(d+2\alpha)(e_2Dv+e_3u^2v^2)
                 +2e_4Dv(e_2+e_3u^2)\big]\\
&+2e_2e_3\big[d(u^2v^2+Dv)+2\alpha+2e_5v^2(1+u^2)\big].
\end{aligned}
\tag{5}
\]
Every term is nonnegative, and d,e2,u,v\textgreater0 give strict positivity.
Thus eta\textgreater0 throughout the contact graph.

\hypertarget{the-three-remaining-sign-walls-cannot-be-crossed}{%
\subsection{4. The three remaining sign walls cannot be crossed}\label{pe31b9ef5ec43:the-three-remaining-sign-walls-cannot-be-crossed}}

At any physical zero, a retained event with zero amplitude can be removed
while preserving the kernel and angular stationarity equations: the event
value and every first derivative of its squared amplitude vanish there.
The stationary-contact theorem makes the resulting nonnegative coefficient
vector physical. A physical d\textgreater0 vector with noninterior Newton support,
however, cannot have a noncommuting zero, by the routing theorem.

If Y=0, the retained e4 event vanishes. Removing it gives omissions
lambda,e1,e4, including the distant pair lambda,e4. If A-Y=0, the retained
e3 event vanishes, and its removal gives the distant omitted pair e1,e3.
Both contradict the preceding paragraph. Therefore Y and A-Y never vanish.

For eta=1, use the exact cross-multiplied contact identity
\[
e_2\,[A+uv(\eta-1)]\,[v(A+u^2Y)-u\eta]
=2u^3(A-Y)(1-\eta)(1+v^2).
\tag{6}
\]
To justify (6) even when a selected Cramer minor vanishes, first derive it
on the nonzero-minor locus from the verified rational contact solution.
The selected minor is an analytic function on the connected positive event
orthant. It is not identically zero, as the exact seed in Section 5 shows.
A nonzero analytic function on a connected open set cannot vanish on an
open subset; hence the nonzero-minor locus is dense. Both sides of (6) are
continuous on the entire analytic contact graph, so (6) extends everywhere.
This explicitly handles exceptional Cramer points.

At eta=1, A\textgreater0 and e2\textgreater0 make (6) force
\(v(A+u^2Y)-u=0\), which is precisely the retained e2 amplitude. Removing
e2 gives the distant omitted pair lambda,e2, again impossible. Hence eta-1
never vanishes.

The continuous nonzero functions Y, A-Y, and eta-1 have fixed signs on the
connected positive event orthant. The exact seed gives their signs.

\hypertarget{exact-positive-seed-and-transfer}{%
\subsection{5. Exact positive seed and transfer}\label{pe31b9ef5ec43:exact-positive-seed-and-transfer}}

Take
\[
u=v=A=1,\quad Y=\tfrac14,\quad\eta=\tfrac12,
\]
\[
c=\left(\tfrac{27}{2},0,0,4,2,9,\tfrac54,
                 \tfrac{13}{4},\tfrac94\right).
\tag{7}
\]
The complete parent kernel and both angle stationarity equations hold
exactly; the selected Cramer minor is nonzero. All six event weights are
positive. Contact sufficiency therefore places it on the graph (2). It
has Y\textgreater0, A-Y\textgreater0, and eta-1\textless0. The previous section proves (1) over the reals.

The final implication is a finite first-order assertion in the coefficient,
state and angle variables: physicality is expressed by quantified
polynomial PSD inequalities; the kernel and stationarity equations are
polynomial; positivity of u,v,d and the six retained events and the desired
signs are ordered-field formulas. Transfer proves (1) over every real
closed extension. No logarithm or analytic continuation is used within a
non-Archimedean field.

\hypertarget{exact-three-consecutive-omissions}{%
\subsection{6. Exact three-consecutive omissions}\label{pe31b9ef5ec43:exact-three-consecutive-omissions}}

After choosing the adjacent omitted pair lambda=e1=0, a three-consecutive
omission has either beta=0 or e5=0. The other five events are positive.
At any real finite noncommuting contact with d\textgreater0, the full Hessian remains
positive definite and the core expectation remains strictly negative by
the same established theorems. Apply the implicit function argument of
Section 1 at this boundary tuple and open its third omitted coefficient
to epsilon\textgreater0, keeping the other event weights fixed. For sufficiently
small epsilon the signed coefficient stays positive and the branch has
nonnegative coefficients. Stationary-contact sufficiency makes it physical.
Thus it belongs to the six-positive-event graph already treated.

Taking the analytic limit gives Y\textgreater=0, A-Y\textgreater=0, and eta\textless=1. Strict A\textgreater0 and
eta\textgreater0 follow directly from Sections 2 and 3; the cofactor formula stays
strict when beta=0 or e5=0. The still-positive e4 and e3 coefficients make
the same distant-omission removal arguments exclude Y=0 and A-Y=0.
Equation (6) passes to the limit along the analytic branch; e2 remains
positive, so eta=1 again forces the forbidden e2 removal. This proves
strict inequalities (1) at the three-omission boundary as well.

This limit is taken at a fixed finite real noncommuting contact, where the
implicit-function hypotheses are verified. It does not assume convergence
of an optimizer escaping to infinity. The resulting first-order implication
then transfers to every real closed extension.

\hypertarget{all-amplitude-signs-and-a-single-rational-chart}{%
\subsection{7. All amplitude signs and a single rational chart}\label{pe31b9ef5ec43:all-amplitude-signs-and-a-single-rational-chart}}

There is a physical Bell reflection preserving the adjacent omitted pair.
It exchanges e2 with e3, e4 with alpha, and e5 with beta (and lambda with
e1). After real local basis gauges its angle action is u'=1/u, v'=v and
its state action, up to a positive normalization, is
\[
\psi'=(vA-u(\eta-1),\ -A-uv(\eta-1),\ 1-uvY,\ uY+v)^T.
\tag{8}
\]
The matrix implementing this state action is
\[
G=\begin{pmatrix}v&-1&0&0\\-1&-v&0&0\\0&0&-v&1\\0&0&1&v\end{pmatrix},
\qquad G^TG=(1+v^2)I.
\]
If c' denotes the displayed coefficient permutation, direct full-parent
congruence gives \(u^2(1+v^2)P_{c'}(1/u,v)=GP_c(u,v)G^T\).
The transformed event weights satisfy the same positivity or single-zero
conditions. Apply (1) to the transformed contact and divide its state by
its positive last coordinate uY+v. The inequalities Y'\textgreater0, eta'\textless1,
eta'\textgreater0, and Y'\textless A' give, respectively,
\[
1-uvY>0,\qquad h_0:=A+uv(\eta-1)>0,
\]
\[
R_5:=u(A-Y)+u^2v(\eta-1)-v<0,
\qquad R_2:=v(A+u^2Y)-u\eta>0.
\tag{9}
\]
Thus every one of the eight event amplitudes has a fixed, strictly nonzero
sign in this canonical state gauge. In particular, the denominator product
in the established selected-minor identity is strictly positive:
\[
\mathcal D=Y(1+u^2)(uY+v)(1-uvY)h_0R_2(-R_5)>0,
\]
\[
\det J_{(d,e_2,e_4,e_5,\alpha,\beta)}
=-8192\eta u^2v(1+u^2)^2(1+v^2)^2\mathcal D<0.
\tag{10}
\]
Consequently the same e3=2 rational Cramer formula is valid on both the
exact two-omission family and its routed three-consecutive-omission
boundaries. It is legitimate here to include a zero coefficient's column
in the selected square matrix: (10), rather than an active-column deletion
argument, proves its invertibility. This replaces the earlier need for six
separate five-row templates on the adjacent three-omission boundary.

For six positive events, (10) also proves that every maximal-signed
boundary point above is an extreme ray of the full canonical cone. The
seven active columns have rank six and their kernel is one-dimensional.
For a single additional event omission, the six active columns have rank
five: the unique kernel of all seven columns has zero in that omitted
coordinate and is strictly positive on the remaining active coordinates.
The exact contact-face description then proves extremality as well.

\hypertarget{what-this-closes}{%
\subsection{8. What this closes}\label{pe31b9ef5ec43:what-this-closes}}

The adjacent exact-two-omission contact formula now enters its positive
shape domain globally: A\textgreater0, t=u/A\textgreater0, 0\textless r=Y/A\textless1, and 0\textless eta\textless1.
The same conclusion holds for the routed three-consecutive-omission contact
families, with the same rational formula and a strictly signed denominator.
Those quantities need not have positive real residues or be bounded.
Consequently this theorem closes sign entry, but not the remaining joint
rate and optimizing-angle boundary certificates. The separated omission
face remains outside this theorem.

\hypertarget{verification-and-review}{%
\subsection{Verification and review}\label{pe31b9ef5ec43:verification-and-review}}

The checker verify\_sign\_entry.py reconstructs the full parent and verifies
the negative-core Schur complement, the positive cofactor identity, the
cross-multiplied contact identity on the Cramer locus, the exact seed kernel,
both angular derivatives, its selected minor -8640, and the positive d=0
central coefficient. The analytic graph, dense extension, no-wall arguments,
and real-closed-field transfer are proved in the text. Two same-system
agents reviewed the argument and its three-omission extension; this is not
external or proof-assistant verification.
